%% sections/02-related-work.tex

\section{Related Work}
\label{sec:related}

\subsection{Playtest Methodology in Game Design}

Playtesting is the primary quality-assurance mechanism in game development, yet it
has received limited systematic treatment in the academic literature. Practitioner
accounts dominate: Fullerton's \textit{Game Design Workshop}~\cite{fullerton2008}
articulates playtesting as iterative discovery of design failures, with evaluators
focused on moments when the game ``breaks'' rather than on assessing the degree to
which designed qualities are present. This deficit-oriented model treats playtesting
as a debugging process. Hunicke, LeBlanc, and Zubek's MDA framework~\cite{playtesting2019}
provides a vocabulary for distinguishing mechanics, dynamics, and aesthetics in game
analysis, but does not provide a scoring instrument.

Zagal et al.~\cite{zagal2008collaborative} study collaborative games and identify
dimensions of cooperative quality (shared decision-making, role complementarity, joint
problem solving) that prefigure several of the Marathon rubric's concerns. Their
analysis is category-level rather than scored, and addresses board games rather than
improvisational narrative games. Deterding et al.~\cite{deterding2011game} characterize
game design elements in terms of their contribution to motivational quality, which
aligns with the Character Agency dimension but does not provide a measurement framework.

The Marathon rubric differs from prior playtest instruments in three ways. First, it
is both deficit-oriented (penalizing failures in Mechanical Fairness, Module Fidelity)
and surplus-oriented (rewarding the presence of emergent positive outcomes in Surprise
and Atmospheric Landing). Second, it produces a numerical score rather than a
categorical judgment, enabling trajectory analysis across sessions. Third, it is
explicitly designed to operate on session logs rather than participant reports,
making it applicable to AI-simulated sessions where no human participants are present.

\subsection{AI as Dungeon Master}

Automated narrative systems have a long research history, from interactive drama
frameworks~\cite{riedl2016} to procedural quest generation~\cite{dormans2010} to
character-driven story engines~\cite{ammanabrolu2021}. The primary evaluation
challenge across this literature is that quality metrics designed for static narrative
(coherence, plausibility, narrative arc completeness) do not capture the improvisational
and participatory qualities specific to role-playing contexts.

Akoury et al.~\cite{akoury2020} introduce STORIUM, a dataset and evaluation platform
for machine-in-the-loop story generation that shares the Marathon corpus's concern for
structured narrative quality evaluation. STORIUM uses human ratings and card-based
constraint systems; it does not address the fully-automated setting where no human
participants are present. Fan, Lewis, and Dauphin~\cite{fan2019} study strategies for
LLM story generation with a focus on narrative structure but evaluate through human
raters reading story outputs, a method that does not scale to large session corpora.

The AI Dungeon system~\cite{aidungeon2019} demonstrates that large language models
can produce coherent narrative in TRPG-like contexts at scale, but its evaluation
relies on player engagement signals (session continuation, sharing behavior) rather
than structural session quality assessment. Mirowski et al.~\cite{mirowski2023}
examine co-writing scripts and theatrical scenes with AI, identifying quality
dimensions (character voice consistency, dramatic arc coherence, staging feasibility)
that partially overlap with the Marathon rubric's Character Agency and Module Fidelity
dimensions. Their evaluation uses theatrical practitioners as raters rather than an
automated scoring instrument.

\subsection{Quality Measurement in Creative and Narrative Systems}

The broader literature on quality measurement for AI-generated text identifies three
primary approaches: reference-based metrics, preference-based metrics, and
property-based metrics~\cite{celikyilmaz2020}. Reference-based metrics (BLEU~\cite{papineni2002bleu},
ROUGE~\cite{lin2004rouge}) measure similarity to a gold-standard corpus; they require
a reference corpus and cannot evaluate properties that have no single correct expression.
Preference-based metrics (pairwise comparison, rating scales) collect human judgments
but do not scale to large corpora without crowdsourcing infrastructure~\cite{akoury2020}.

Property-based metrics check that outputs satisfy specified properties, independently
of any reference. Andrade~\cite{andrade2000} develops scoring rubrics in educational
assessment contexts that are explicitly property-based: each rubric criterion defines
a property that a successful work should exhibit, and the rubric score reflects the
degree to which evidence of that property is present. This is the model the Marathon
rubric follows.

LLM-as-judge approaches~\cite{zheng2023judging} have gained prominence as a scalable
alternative to human rater studies, using a language model to evaluate outputs from a
language model. Liang et al.~\cite{liang2023} apply this approach to research paper
review with encouraging results for factual dimensions. The Marathon rubric explicitly
avoids this approach for its primary scoring, opting for rule-based evidence citation
(a specific passage from the log must be cited for every non-zero sub-score) rather
than model judgment. However, the five-lens panel review that runs alongside the rubric
gate does use AI-as-expert simulation, following the approach described by Gero and
Kan~\cite{gero2004} for design critique and extended by Mirowski et al.~\cite{mirowski2023}
for creative work evaluation.

Yuan et al.~\cite{yuan2022wordcraft} study collaborative story writing with LLMs using
a tool that offers completions, continuations, and elaborations to human authors.
Their evaluation focuses on the human author's experience (satisfaction, perceived
quality, creative ownership), which is orthogonal to the Marathon rubric's focus on
automated session log quality. Lee, Liang, and Yang~\cite{lee2022coauthor} similarly
study human-AI co-writing and find that the dimensions humans care about (agency,
expressiveness, surprise) are multi-dimensional and do not reduce to a single
quality axis, a finding consistent with the Marathon rubric's 8-dimensional structure.

\subsection{Expert Review Simulation}

Simulated expert review --- using an AI system to adopt the perspective of a domain
expert and evaluate creative or technical work --- has been studied in design
critique~\cite{gero2004}, literary and dramatic evaluation~\cite{mirowski2023}, and
research peer review~\cite{liang2023}. The common finding is that structured role
assignment (giving the AI a specific expert identity and evaluation rubric) significantly
improves the specificity and consistency of the produced evaluation compared to open-ended
critique requests.

The Marathon workshop uses two expert-simulation panels operating at different points
in the session pipeline. A five-persona design review panel (Gygax, Jaquays, Schick,
Moldvay, Patrick Stuart) evaluates compiled adventure modules before play, scoring on
axes specific to each designer's documented aesthetic priorities. A five-lens
player-experience panel (Tactician, Roleplayer, Lorekeeper, Improviser, Fresh Face)
evaluates session logs after play from five distinct player perspectives. These panels
are distinct from the rubric-based session gate and serve complementary rather than
redundant functions: the gate scores evidence of outcomes in the session log; the panels
evaluate whether those outcomes would satisfy specific types of players or design critics.

Tychsen et al.~\cite{tychsen2006role} study how player type affects TRPG experience
across platforms, identifying that the same session can succeed for one player type and
fail for another. This finding motivates the Marathon rubric's approach of using multiple
player-experience lenses in the panel rather than a single aggregate player judgment.
